
\capitulo{Proceso Verificación}
\label{sec:srs-5}

El proceso de verificación comprobará que el producto satisfaga los requerimientos definidos en la sección 3, mediante revisiones, inspecciones, demostraciones y pruebas durante las iteraciones de desarrollo. Se tomarán como referencia el estándar de procesos del ciclo de vida del software~\cite{iso12207_2026} y la serie de estándares de pruebas de software~\cite{iso29119_2021}, de acuerdo con los procedimientos y criterios de medición establecidos en esta especificación.

La planificación se basará en la tabla de referencia cruzada de la sección 2.7, que relaciona las funcionalidades, los módulos, los casos de uso y los requerimientos funcionales. Esta trazabilidad permitirá asociar los requerimientos funcionales RF-01 a RF-14 y los no funcionales RNF-01 a RNF-14 con las verificaciones correspondientes, considerando también los requerimientos de interfaces, desempeño y base de datos de la sección 3.

Las pruebas funcionales comprobarán la consulta y corrección de condiciones meteorológicas (RF-01 a RF-03), la ubicación del incendio y recuperación de escenarios (RF-04 a RF-08), la agregación de resultados (RF-09 a RF-11) y su visualización (RF-12 a RF-14). Las pruebas no funcionales evaluarán la confiabilidad (RNF-01 a RNF-05), disponibilidad (RNF-06 y RNF-07), mantenibilidad (RNF-08 a RNF-11) y portabilidad (RNF-12 a RNF-14), junto con los límites de desempeño establecidos en la sección 3.3.

Los casos de prueba se diseñarán a partir de los criterios de medición de cada requerimiento, utilizando las técnicas de la parte 4 de la serie citada: partición de equivalencia, para comprobar entradas meteorológicas válidas e inválidas (RF-02); análisis de valores límite, para evaluar los extremos permitidos de las variables meteorológicas y los tiempos de espera; tablas de decisión, para verificar combinaciones de disponibilidad de Open-Meteo y compatibilidad de escenarios (RF-01, RF-02 y RF-07); y pruebas de transición de estados, para comprobar la actualización de consultas y el comportamiento ante archivos faltantes o dañados (RF-08 y RNF-05).

La ejecución de las pruebas se documentará conforme a la parte 3 de la misma serie, registrando el identificador del requerimiento, las condiciones de prueba, el resultado esperado, el observado y las evidencias. Los incumplimientos se registrarán y corregirán, repitiendo las pruebas afectadas y ejecutando pruebas de regresión cuando sea necesario.

La tabla de referencia cruzada permitirá dar seguimiento a los requerimientos verificados, pendientes o incumplidos. Los resultados complementarán las actividades de validación con usuarios descritas en la sección~\ref{sec:srs-4-4}, sin sustituir la evaluación de utilidad y usabilidad prevista con los bomberos.
